\documentclass{article}
%\usepackage{corl_2023} % Use this for the initial submission.
 \usepackage[final]{corl_2023} % Uncomment for the camera-ready ``final'' version.
%\usepackage[preprint]{corl_2023} % Uncomment for pre-prints (e.g., arxiv); This is like ``final'', but will remove the CORL footnote.
\usepackage{varwidth}
\usepackage{multirow}
\usepackage{enumitem}
\usepackage{booktabs}
\usepackage{threeparttable}
\usepackage{multirow}
\usepackage{times}
\usepackage{epsfig}
\usepackage{graphicx}
\usepackage{amsmath}
\usepackage{amssymb}
\usepackage{algpseudocode}
\usepackage[utf8]{inputenc}
\usepackage[longend]{algorithm2e}
\usepackage{textgreek}
\usepackage[font=small]{caption}
\usepackage{pdfpages}
\usepackage{adjustbox}
\usepackage{wrapfig}
\usepackage{array}
\usepackage{sidecap}


\usepackage{enumitem}
\title{SIMF: Semantics-aware Interactive Motion Forecasting for Autonomous Driving}

%Trajectron++: Dynamically-Feasible Trajectory Forecasting With Heterogeneous Data

%Scene Semantics-aware interactive motion forecasting for autonomous driving

% The \author macro works with any number of authors. There are two
% commands used to separate the names and addresses of multiple
% authors: \And and \AND.
%
% Using \And between authors leaves it to LaTeX to determine where to
% break the lines. Using \AND forces a line break at that point. So,
% if LaTeX puts 3 of 4 authors names on the first line, and the last
% on the second line, try using \AND instead of \And before the third
% author name.

% NOTE: authors will be visible only in the camera-ready and preprint versions (i.e., when using the option 'final' or 'preprint'). 
% 	For the initial submission the authors will be anonymized.

\author{
  Vidyaa Krishnan Nivash, Ahmed H. Qureshi\\
  %Department of Electrical and Computer Engineering\\
Purdue University\\ 
  West Lafayette, IN 47907, United States\\
  \texttt{\{krish255, ahqureshi\}@purdue.edu} \\
  %% \AND
  %% Coauthor \\
  %% Affiliation \\
  %% Address \\
  %% \texttt{email} \\
  %% \And
  %% Coauthor \\
  %% Affiliation \\
  %% Address \\
  %% \texttt{email} \\
  %% \And
  %% Coauthor \\
  %% Affiliation \\
  %% Address \\
  %% \texttt{email} \\
}


\begin{document}
\maketitle

%===============================================================================
\vspace{-10pt}
\begin{abstract}
Autonomous vehicles require motion forecasting of their surrounding multi-agents (pedestrians and vehicles) to make optimal decisions for navigation. The existing methods focus on techniques to utilize the positions and velocities of these agents and fail to capture semantic information from the scene. Moreover, to mitigate the increase in computational complexity associated with the number of agents in the scene, some works leverage Euclidean distance to prune far-away agents. However, distance-based metric alone is insufficient to select relevant agents and accurately perform their predictions. To resolve these issues, we propose Semantics-aware Interactive Motion Forecasting (SIMF) method to capture semantics along with spatial information, and optimally select relevant agents for motion prediction. Specifically, we achieve this by implementing a semantic-aware selection of relevant agents from the scene and passing them through an attention mechanism to extract global encodings. These encodings along with agents' local information are passed through an encoder to obtain time-dependent latent variables for a motion policy predicting the future trajectories.  Our results show that the proposed approach outperforms state-of-the-art baselines and provides more accurate predictions in a scene-consistent manner. 
    %Autonomous vehicles require motion forecasting to find an optimal speed to reach their destination without any perturbations to existing traffic. It helps them to act proactively, when changing lanes, overtaking other participants, managing intersections, etc. Forecasting multi-agent motion can be inspired by human intuition, boiling down to two broad aspects of human visual perception: (1) ambient mode, based on the positions and velocities of the participants, and (2) focal mode, based on identifying the nature of the participants. Many of the prior works are based only on the ambient mode or with minimal, implicit usage of the focal mode. Also, they lack the active interaction between these two modes to extract the complete information of the rich multi-modal data captured from the surroundings. In this paper, we present SIMF, which simultaneously extracts information from both modes and allows active interaction between them. Specifically, we achieve this by (1) Using semantics from map encodings to explicitly capture the focal mode information, (2) Capturing global contextual information between the semantics and spatial orientation by passing them through an attention mechanism, and (3) leveraging a time-dependent latent encoding technique to capture dynamics from both the modes of visual perception. 
\end{abstract}
% Two or three meaningful keywords should be added here
\keywords{autonomous driving, scene semantics, multi-agent motion forecasting} 

%===============================================================================

\section{Introduction}\vspace{-10pt}
Driving is a deceptively simple activity requiring continuous visual sampling, periodic control, and sustained attention \cite{CASTRO}. Humans simplify this task by sampling the environment for critical control cues, typically involving predictable extrapolations of well-learned circumstances. Human cognition can perform these perceptual estimations with high accuracy to ensure safe driving. According to \cite{Tyrrell}, two separate visual perceptual mechanisms exist to make good estimations. The first being the ambient mode, concentrates on factors such as spatial orientation and locomotion. It involves vehicle guidance using information gained from optical flow to derive estimations of vehicle positioning. The second being the focal mode, focuses on object recognition, identification, and the semantics associated with it. It concentrates on the detection and identification of objects of importance in the environment such as physical barriers, and roadway threats, and uses its practical semantic associations as critical cues for deriving estimations. 
 
%In order to achieve considerable accuracy in replacing  humans with autonomous vehicles, these estimations have been studied in greater detail. The current research on human driving behavior prediction has been developing methods that learn the behavior of agents from large data i.e. from the sensors and vehicle inputs. Yet, most of the current work focuses explicitly on the ambient mode and takes only implicit cues from the second by inherently learning from data. %Specifically, it learns a policy from the agents' history of states containing its position, orientation, and velocity information irrespective of knowing each agent's relevance to the ego vehicle in terms of its semantics associated with the scene. 
 %In this paper, we present SIMF, an extensible approach built upon the ScePT framework to introduce the focal mode and imbibe an active interaction between the two modes of perceptual mechanisms. Our key contributions are three-fold: Firstly, We introduce techniques to optimize the grouping of relevant agents in a scene by formulating an algorithm to extract the semantics of agents and scene elements from map semantics. Secondly, we combine both the modes of perceptual mechanisms, by feeding the local map and the state and edge history information through an attention mechanism to obtain global encodings.
%Specifically,  we remove redundant cliques and capture temporal dependencies in the clique by introducing \textbf{a sliding window approach throughout the timestamps }to capture the time information in agents. 
%Thirdly,  we capture inter-clique interactions and new relevances in the data with respect to time using a method to encode temporal information. %and capture implicit interrelations among all the agents. 


The current work in designing autonomous driving motion prediction methods focuses explicitly on the ambient mode and takes only implicit cues from the second, i.e., the focal mode, by inherently learning from data. Hence, the lack of context from the explicit modeling of the focal mode and its inherent interaction with the ambient mode limits the existing methods' performance accuracy and efficiency in complex environments. Therefore, in this paper, we present SIMF, an extensible approach to introduce the focal mode and imbibe an active interaction between the two modes of perceptual mechanisms. Our key contributions are three-fold: Firstly, we formalize focal mode by formulating an algorithm to extract the semantics of agents and scene elements from environment maps. This extracted information is then utilized to optimize the grouping of relevant agents in a scene for their scene-consistent motion forecasting. Secondly, we combine both modes of perceptual systems through an attention mechanism to obtain their global, interactive representation. Thirdly,  we capture inter-group agents' interactions and new relevances in the data in a time-series manner using temporal encodings of focal and ambient representations. Finally, the motion prediction policy uses these temporal encodings to predict agents' future trajectories under scene consistency. Hence, SIMF aims to inculcate an active interaction between the two visual systems and induce an ability to complement one another, achieving mutual reinforcement for trajectory predictions. This results in an overall reduction in response time and improves the accuracy of trajectory predictions compared to state-of-the-art (SOTA) baseline methods. 

\section{Related work}\vspace{-10pt}
Research in trajectory prediction has explored multiple facets of driving, uncovering a broad spectrum of approaches and possibilities.  Early works majorly delved into decision-making processes based on deterministic models like Dynamic Bayesian Networks \cite{Lefvre2014ASO}, Gaussian
 Process Regression (GPR) \cite{4359316} and Recurrent Neural Networks (RNNs) \cite{article} \cite{8460504} \cite{jia2022towards} . According to the review papers \cite{Leon_2021} \cite{Rudenko_2020}, motion prediction techniques can be classified based on the number of modes, level or type of abstraction, and method of approach. 
%The comprehensive classification of human motion prediction is fully depicted in this survey paper \cite{Rudenko_2020}.
Many data-based approaches in the literature rely on LSTMs
%Long Short-Term Memory (LSTMs) 
or simpler CNNs
%convolutional neural networks (CNNs)
\cite{nikhil2018convolutional} to learn from agent states and histories. But these techniques do not ensure scene consistency, i.e., the trajectory predictions of multiple agents collide with each other or with static obstacles. 
%This issue arises due to independent predictions for each timestep of individual agents and the lack of ability to capture dynamic implicit information.  
To resolve this issue, recent studies have employed conditional variational autoencoders (CVAEs) or more recent methods
such as generative adversarial networks (GANs) and attention mechanisms to achieve SOTA work in motion prediction for autonomous vehicles. 

%Higher abstraction levels in prediction models tend to increase computational complexity, so algorithms have been proposed to focus on the most plausible trajectories. Risk estimation is also crucial for prediction methods, as highlighted in \cite{Lefvre2014ASO}, which emphasizes the importance of considering both vehicle motion modeling and risk estimation. 
% In [7], another LSTM model is described for interaction-aware motion prediction.
% Confidence values are assigned to the maneuvers that are performed by vehicles. Based on
% them, a multi-modal distribution over future motions is computed. More specifically, the
% model computes probabilities for each type of maneuver, based on six maneuver classes.
% The input to the LSTM is represented by the past positions of the ego car and its neighbors,
% and the geometry of the road lanes.
% Social LSTM [8], used for predicting the trajectory of pedestrians, uses LSTM with
% a social pooling layer which allows neighbors, up to a certain distance, to exchange
% information. The hidden states of their corresponding LSTMs are pooled together and used
% as an input for the following prediction step. \\ \\ 


%\cite{girgis2022latent} uses Sequential Set Transformers for producing latent variables with temporal information for joint multi-agent motion predictions. Nonetheless, when applied to complex scenes involving multiple agents, transformers incur high computational time and complexity. In other words, the computational complexity grows exponentially as the number of agents increases.

\begin{wraptable}{r}{5.5cm}\vspace{-0.18in}
\begin{center}
 \setlength\tabcolsep{-0.75pt}
 \vspace{-8pt}
\begin{tabular}{cccccc}

\hline Method & $\mathrm{SC}$ & $\mathrm{TE}$ & $\mathrm{AF}$ & $\mathrm{SI}$ & $\mathrm{OC}$  \\
\hline $\begin{array}{l}\text {SRNN}[13]
\\
\text {SSTran}[12] \\
\text {LaFormer}[15]\\
\text {LaneGraph}[14] \\
\text {ScePT}[16] \\
\text{SIMF}\end{array}$ & $\begin{array}{l}\\
\checkmark \\
\\
\\
\checkmark \\
\checkmark
\end{array}$ & 
$\begin{array}{l}
     \\ 
    \checkmark \\
    \checkmark \\
    \checkmark \\
    \\
    \checkmark
     
\end{array}$
& $\begin{array}{l}
\\
 \\
 \\
\\
\\
\checkmark
\end{array} $
& $\begin{array}{cc}
   \\
   \\
   \checkmark\\
  \checkmark \\
   \\
   \checkmark
\end{array}$
& $\begin{array}{l}
\checkmark \\ \\
\\
\\
\checkmark
\\
\checkmark\end{array}$\\
\hline

\end{tabular}
\end{center}
\caption{Legend: SC - Scene Consistency, TE - Temporal Encoding, AF - Ambient and focal mode interaction, SI - Semantic-aware Interaction, OC - Optimal computation time}
\vspace{-0.2in} \label{table1}\end{wraptable}
However, these techniques have time-independent latent variables and hence do not capture temporal dependencies within each agent's history and dynamic relationships between agents. For example, in CVAE, the latent variable encodings for the input agent states and histories are constant with respect to time. Hence they do not explicitly learn the temporal graph governing the causal relationships in its inputs. The paper \cite{7796911} \cite{girgis2022latent} state this problem for general time-series data and proposes temporally dynamic latent variables to account for the lag in its memory states.  Some authors tend to resolve this issue using RNN for encoding the temporal information \cite{jain2016structuralrnn} \cite{jia2022towards}, but these methods cannot handle multimodal trajectory predictions. Similarly, \cite{lanegraph} predicts trajectories based on lane graphs by exploring both lateral (e.g.: lane keeping, turning) and longitudinal (e.g.: acceleration, braking) uncertainties. However, they do not capture the interaction between the two modes. Also, such works \cite{laformer} focus on lane semantics, thus the other semantics (roadblocks, pedestrian walkways, barriers) are left unexplored.
%AgentFormer %\cite{agentformer} proposes to capture the time and social dimension by leveraging a sequence representation of multi-agent trajectories by using an attention mechanism. But they do not account for the increase in computational complexity while working with complex and dense scenes. 

Another recent and perhaps the most relevant method to our approach is ScePT \cite{chen2022scept} framework, which follows an ambient mode of perception, ensuring scene consistency and computational efficiency. It is a discrete CVAE model that outputs joint trajectory predictions and achieves scene consistency by using an explicit collision penalty as regularization during training. Furthermore, it reduces the computational complexity by forming subgraphs of agents, also known as cliques, based on Euclidean distance and jointly predicts their future states rather than for an individual agent or a whole set of agents. Although ScePT provides SOTA performance, it lacks the following features that limit its performance. First, it does not account for the temporal dependencies of the agents, which, as highlighted earlier, is crucial for better performance. Second, it only considers the ambient mode of perception and does not have a focal mode based on agents' local perception fields. Lastly, it forms cliques based on Euclidean distance only and does not account for other scene semantics such as lanes, road barriers, etc. Therefore, in this paper, we propose SIMF, a full-stack motion prediction model that, with ambient and focal modes of perception, considers all possible scene semantics and temporal dependencies of selected agents, as summarized in Table ~\ref{table1}, and provides better performance among all baselines.
%One method that allows to attain all these objectives {\color{red} which objectives?, specify} is the ScePT \cite{chen2022scept}, a discrete CVAE model that outputs joint trajectory
%predictions for multiple agents in a scene, ensuring high scene consistency in its predictions. It ensures high scene consistency by using an explicit collision penalty as regularization during training. It also reduces the computational complexity by using a novel technique {\color{red} Which technique? do you mean using a novel Euclidean distance based clique, subgraph, formation ?} on a CVAE model to reduce the dimensionality of the product latent space in the encoder of CVAE architecture. In order to improve ScePT, we include aspects of the temporal encodings and focal mode of visual perception using techniques mentioned in Section 4. The table provides a detailed breakdown of recent SOTA approaches and their consideration of various factors.
% Taking into account the time constraints of a real-time system, [9] uses a
% simple feed-forward CNN architecture for the prediction task. The authors use an RGB
% image to represent the scene context. However, a vector of velocity, acceleration, and yaw
% rate can also be included. In this case, this vector is concatenated with the flattened output
% of the CNN. Then, these aggregated features are sent to a fully connected layer.
% A similar approach is used in [10], which predicts multiple possible trajectories together
% with their probabilities. The context is also encoded as an image that is passed to a
% CNN. Given the raster image and the state estimates of agents at a time step, the CNN is
% used to predict a multitude of possible future state sequences, as well as the probability of
% each sequence. Overall, most of the existing works involve focusing on the ambient mode of the visual system and completely ignore or do not use the focal mode explicitly. Our work introduces a new technique for combining the modes and capturing the interaction between these two modes thereby increasing accuracy and reducing response time. 
%\section{Background}  
%Our work presents a method to generate joint trajectory predictions of all the interacting agents in a scene. These agents are captured as a spatiotemporal graph $G = (V,E)$ from the scene, where $V$ represents the nodes and $E$ represents the graph's edges. The agents are represented as nodes, and the interactions between them are represented by the edges. Henceforth, we use the terms nodes and agents interchangeably. 
%These agents are broadly classified into vehicles and pedestrians, utilized throughout the model to determine the dynamics $\psi$ of each agent based on its respective semantics. The state of a node is represented by $s\ \epsilon\ \mathbb{R}^4$ containing information about the position $p = (x,y)$, orientation ($\theta$), and velocity ($v_x$, $v_y$). Specifically, $s_v$ contains vehicle state information $s_v = (x,y, v, \theta)$,  where $v$ is the $L_2$ norm of $(v_x, v_y)$ and $s_p$ contains pedestrian information $s_p= (x,y,v_x,v_y)$. The states are stacked together from past $H$ timesteps to timestep $t$, given by $\vec h^t_s = (s^{t-H}, s^{t+1-H}, ..., s^t)$ while $\vec h^t_e$ represents the corresponding stack of edge information. We use the state information to model an edge between a pair of nodes. The edge quantifies the degree of interaction ($\alpha_{ij} \in [0,1]$), where $(i,j)$ represents the pair of nodes connected by an edge.  
%$\alpha_{ij}$ for each edge is measured by a series of steps mentioned in Algorithm 1.  \\ 
%Our work assumes that maps containing geometric and semantic information are available. These maps are converted into local maps around each node. With a particular pixel resolution $r$, we consider a  fixed $K \times K$ size for all local maps, each modified into having $L$ semantic channels given by $M^t \ \epsilon \ \mathbb{R}^{K \times K \times L} $. The map semantic channels are constituted by the token identifiers $\tau$. For example, roadblocks are defined by a set of edge line tokens $\tau_{e_{0:\lambda}}$ where $\lambda$ represents the total number of tokens required to define the spline. Similarly, lanes are represented by curve spline tokens $\tau_l$. We form cliques $C^t$, i.e., subgraphs of agents having a high $\alpha_{ij}$ for a set of sampled timesteps $t \  \epsilon \  (0,T)$, $T$ representing the total number of timestamps considered for clique formation. We form the cliques based on the state and edge histories along with map semantics given by $ C^t_{1:n} = \sigma(s_{t_{1:N}},\vec h^t_s, \vec h^t_e, M^t )$. The $n$ represents the number of cliques formed from $N$ number of agents in the scene. The function $\sigma$ leverages various road semantics, $M^t$, at time $t$, such as lane barriers, crosswalks, lane directions, the distance between agents, etc., to quantify the edge value $\alpha$ between nodes and forms cliques of nodes with $\alpha$ greater than a threshold. Finally, given a clique, $C^{t:t-H}$, agents' states, $h^t_s$, and interactions, $h^t_e$, at time $t$, we aim for a trajectory prediction function $f$ that predicts the future trajectories of each agent $s_{pred}^{t:t+\eta}$ for a horizon length $\eta$ based on combined measures of the focal and ambient mode of visual perception.$$s_{pred}^{t:t+\eta}\gets f(h_s^{t}, h_e^{t}, C^{t:t-H})$$In contrast to our objective, ScePT \cite{chen2022scept} addresses this problem by using techniques inspired by the ambient mode of visual perceptual mechanism, and does not account for temporal dependencies in the latent variable encodings.  It predicts the future of cliques of agents $C_{t_{1:n}}$, rather than individual agent predictions or the scene graph $G$ as a whole. These cliques are derived by using the Euclidean distance parameter $d_{ij}$ as the function $P$ for estimating $\alpha_{ij}$ (Section 4.1). It produces the target distribution by representing the input observations as latent variables constant with respect to time and these variables are fed into a decoder having a policy network that autoregressively rolls out closed-loop trajectory predictions by mapping them to control inputs.  \\  
%We can infer that the $d_{ij}$ metric used for calculating $\alpha_{ij}$ does not model the interaction between agents accurately. Since the entire joint trajectory predictions fundamentally depend on the quality of the cliques, it is highly essential that the cliques be optimal. Also, the latent variable encodings are constant with respect to time and hence do not capture dynamic information explicitly. In addition to that, the map encodings consist of only the local contextual information and miss out on the correlation between the nodes in terms of the agents' semantics. In other words, the encoding mechanisms are insufficient to infer the global semantics and its corresponding associations with its spatial orientations and locomotions.  In short, the current method focuses more on the ambient mode, lacking the aspects of bringing together both ambient and focal modes. Moreover, it fails to capture a correlation between agents' semantics i.e. the global context. This can affect scenarios where there is an issue of occlusion handling in lanes or a lack of cross-traffic alert sensors. 
 % \subsection{Notation}
% In our paper, we use $VEH$ to refer to vehicles and $PED$ for pedestrians. The state of a node is represented by "s", while "$h_s$" and "$h_e$" denote the state history of each node and its corresponding edge history, respectively, from timestep 0 up to a maximum timestep "T". Each node contains state information about the position ($x$,$y$), orientation ($\theta$), and velocity ($v_x$, $v_y$) in both directions  Since our encoder-decoder model is a CVAE, we use z to denote the hidden latent variable encoding. Subscripts i and j indicate the node pointers, where a single letter denotes an individual node, and a pair of letters denotes a pair of nodes.
% \subsection{Contributions}

% In this paper we present (dynamic ScePT), an extensible approach built upon ScePT framework to introduce the focal mode of data perception and imbibe an active interaction between the two modes of perceptual mechanisms. Our key contributions are three-fold: Firstly, We introduce techniques to optimize the grouping of relevant agents in a scene by formulating an algorithm $P(s_{t_{1:N}},\vec h^t_s, \vec h^t_e, \sigma_{t_{1:N}} )$ by extracting the semantics of agents and scene elements from $\sigma$ (Algorithm 1). Secondly, we capture inter-clique interactions and new relevances in the data with respect to time by pruning the cliques iteratively with respect to its active timestamps. Specifically,  we remove redundant cliques and introduce nodes of greater relevance in a clique by having a sliding window approach throughout the timestamps to capture the dynamic information in agents. Thirdly, we combine both the modes of perceptual mechanisms by passing the local map encodings and the state and edge history information $<\vec h^t_s, \vec h^t_e, \sigma_{t_{1:N}}>$  through an attention mechanism to obtain global encoding to obtain implicit interrelations among all the agents.
\section{Proposed Method}\vspace{-10pt}
Our work presents a method to generate joint trajectory predictions of all the interacting agents in a scene. These agents are captured as a spatiotemporal graph $G^t = (V,E)$ for each timestep $t$ from the scene, where $V$ represents the nodes and $E$ represents the graph's edges. The agents are represented as nodes, and the interactions between them are represented by the edges. Henceforth, the terms nodes and agents are used interchangeably,  and are broadly classified into two types, namely vehicles and pedestrians utilized throughout the model to determine the dynamics $\psi$ of each agent based on its respective semantics. The state of a node is represented by $s\ \epsilon\ \mathbb{R}^4$ containing information about the position $p = (x,y)$, orientation $\theta$, and velocity ($v_x$, $v_y$). Specifically, $s_v$ contains vehicle state information $s_v = (x,y, \nu, \theta)$,  where $\nu$ is the $L_2$ norm of $(v_x, v_y)$ and $s_p$ contains pedestrian information $s_p= (x,y,v_x,v_y)$. The states are stacked together from past $H$ timesteps to timestep $t$, given by $\vec h^t_s = (s^{t-H}, s^{t+1-H}, ..., s^t)$ while $\vec h^t_e$ represents the corresponding stack of edge information. We use the state information to model an edge between a pair of nodes. The edge quantifies the degree of interaction ($\alpha_{ij} \in [0,1]$), where $(i,j)$ denotes the pair of nodes.  
%$\alpha_{ij}$ for each edge is measured by a series of steps mentioned in Algorithm 1.  \\ 

Our work assumes that maps containing geometric and semantic information are available. These maps are converted into local maps around each node. For a particular pixel resolution, we consider a corresponding fixed $K \times K$ size for all local maps, each modified into having $L$ semantic channels given by $M^t \ \epsilon \ \mathbb{R}^{K \times K \times L} $. The map semantic channels are constituted by the token identifiers $\tau$. For example, roadblocks are defined by a set of edge line tokens $\tau_{e_{0:N}}$ where $N$ represents the total number of tokens required to define the spline. Similarly, lanes are represented by curve spline tokens $\tau_l$. We form cliques $C^t$, i.e., subgraphs of agents having a high $\alpha_{ij}$ for a set of sampled timesteps $t \  \epsilon \  (0,T)$, $T$ representing the total number of timestamps considered for clique formation. We form the cliques based on the state and edge histories 
 along with map semantics given by $ C^t_{1:n} = \sigma(s_{t_{1:N}},\vec h^t_s, \vec h^t_e, M^t )$, where $n$ represents the number of cliques formed from $N$ number of agents in the scene. The function $\sigma$ leverages various road semantics from $M^t$, at time $t$, such as lane barriers, crosswalks, lane directions, the distance between agents, etc., to quantify the edge value $\alpha$ between nodes and forms cliques of nodes with $\alpha$ greater than a threshold. Finally, given a clique, $C^{t:t-H}$, agents' states, $h^t_s$, and interactions, $h^t_e$, at time $t$, we aim for a trajectory prediction function that predicts the future trajectories of each agent $s_{pred}^{t:t+\eta}$ for a horizon length $\eta$ based on combined measures of the focal and ambient mode of visual perception.
 \begin{figure}
\begin{center}

\includegraphics[width=1.0\linewidth]
   {semantic.pdf}
  
\end{center}
   \caption{ Shows a step-by-step illustration of how the cliques are pruned using semantic information in order to capture the focal mode of visual perception. The orange vehicle represents the ego-vehicle ($A$), differently colored circles represent pedestrians, and the cars represent other vehicle agents. Fig ~\ref{fig:clique}(1) shows the direction criterion effectively pruning agents which move away from the ego-vehicle. Fig ~\ref{fig:clique}(2) shows how cliques are formed based on a distance threshold from $A$.  Fig  ~\ref{fig:clique}(3) represents the barrier criterion pruning agent $B$ based on road semantics. Fig ~\ref{fig:clique}(4) illustrates how the lane semantics affect the clique formation.}
    \label{fig:clique}
\end{figure}
\subsection{Semantics-aware Clique Formation}
The task of predicting the trajectories of all agents in a scene  without collisions is complicated due to factors such as varying entry and exit timestamps and the interdependence between the multi-modal trajectories of different agents. To ensure scene consistency, each agent should be a part of the joint latent distribution, than individual predictions independent of its neighbors i.e. each agent should be associated with a discrete latent variable with cardinality $N$. Consequently, the cardinality $N$ grows exponentially with the number of nodes in the scene. To solve this increase in complexity, the scene can be broken down into smaller groups, called cliques formed based on degree of interaction. Fig \ref{fig:clique} shows techniques to form cliques based on various criteria, such as the distance, direction, and semantics between the agents. These cliques are optimized and pruned to include only the most relevant agents for decision-making, further explained in supplementary material (Algorithm 1).% while eliminating cliques that contain redundant information.
\subsubsection{Euclidean distance parameter} 
Given a scene with multiple nodes, a spatio-temporal scene graph is generated where nodes represent agents and edges represent the interactions between these agents \cite{chen2022scept}. Each agent’s closest future distance is ascertained as a measure of interaction, propagating forward each node $i$, according to an action $a_i$ with a constant velocity model $\phi_{a_i}^t$. The $\phi_{a_i}^t$ maps the initial state of agent, $i$, to $t$ time steps ahead into the future states based on $a_i$ executed at the initial state, at time $t=0$. %until T is the prediction horizon. 
The closest future distance between two agents is defined as $d_{ij} = \min ||\phi_{a_i}^t(s_i), \phi_{a_j}^t(s_j)||_2$. We then define $\alpha_{ij}$ to populate the scene graph adjacency matrix as follows: the $\alpha_{ij} =0$ when $d_{ij} > d_0$ and $\alpha_{ij} =\frac{d_0}{d_{ij}}$ when $d_{ij} < d_0$, where $d_0$ is the threshold distance fixed according to the agent type. 
%\begin{equation}
%  \alpha_{ij} =
%    \begin{cases}
%      0 & d_{ij} > d_0\\
%      \frac{d_0}{d_{ij}} & d_{ij} < d_0\\
     
%    \end{cases}       
%\end{equation}
%where $d_0$ is the threshold distance fixed according to the agent type. 
\subsubsection{Interaction detection using direction}

To ensure that the cliques do not include agents with minimal or no interaction, we categorize the agents based on whether they are moving toward or away from each other.\begin{wrapfigure}{r}
{0.278\textwidth}\vspace{-10pt}\includegraphics[width=3.8cm,height=2.78cm]{interaction.pdf}
\caption{Estimating relative interaction between two agents in a scene.}\vspace{-20pt}
\label{fig:dirn}
\end{wrapfigure}For example, even if agents are within a distance threshold $d_0$, they will have no impact on each other if they are moving in opposite or perpendicular directions. One way to compute this is by taking the ego agent's viewpoint as the observer having position and velocity $(p_o, v_o)$ and determining $(p_i,v_i)$ of the other agents relative to the observer's frame of reference, as illustrated in Figure ~\ref{fig:dirn}. 
The interaction between the ego vehicle and any other agent is considered relevant when there is a timestep within $\eta$ such that the two velocity vectors intersect or lie within a relative distance threshold $D$ \cite{HJM2010ICAS}. Formally, 
$(p_o + t v_o) - (p_i + t v_i) = \vec 0 {\iff} V_o \cap V_i$. Since, $ (p_o + t v_o) - (p_i + t v_i) = (p_o - p_i) + t (v_o - v_i),$ 
we check if the $||(p_o - p_i) + t (v_o - v_i)||_2 < D$. This technique also implicitly ensures scene consistency by introducing direction-based collision checking during clique formation. 

% where the first inequality checks if the two agents towards each other are within a distance threshold, and the second equality checks if they are towards each other. This implies we can directly check if $$ D \ge (s_j - s_i) \times (v_j - v_i) \ge 0$$ 
\subsubsection{Barrier detection}
%include focal mode info
As mentioned earlier, the focal mode of visual cognition involves understanding the meaning and significance of the perceived elements in the environment as the first step for determining the next action. Specifically, humans imbibe critical cues from the perceived elements based on their knowledge, past experiences, and language understanding. For example, according to human cognition, it is trivial to conclude that two agents can not actively interact with each other if there are physical barriers present between them. But this intuitive fact is based on semantic information about the physical barriers. Hence, we implement this technique to detect barriers between agents to avoid clique formation between completely irrelevant agents. 
To understand the formation of cliques better, let us consider the scene in Figure ~\ref{fig:clique}(3).
In this scene, agents A and B are moving toward each other, but there is no possible interaction between the agents due to the physical barrier (road divider) between them. Hence these agents need not be considered together  in the same clique.  To check if there are any barriers between the two agents, map semantics can be used which are derived from scene segmentation. Given these map attributes as line segments dividing roads, we check if there are dividers between any two agents. Consider a line segment $\tau_{A \rightarrow{B}}$ drawn between the agents $A$ and $B$,$\tau_{C \rightarrow D}$ drawn between $C$ and $D$ in Figure ~\ref{fig:clique}(3) respectively. The line segments $\tau_{A \rightarrow{B}}$ and $\tau_{C\rightarrow D}$ intersect if $ACD$ and $BCD$ have opposite directions meaning either $ACD$ or $BCD$ is counterclockwise but not both. To check the direction of $ACD$, we calculate the slopes $m_{AC}$ and $m_{AD}$ of $AC$ and $AD$, respectively.  If $m_{AC} < m_{AD}$, then $ACD$ has a clockwise direction, else otherwise. We use this condition to check %the directions of $ACD$ and $BCD$ and if they are having opposite direction of rotations, then we infer that the agents are separated by physical barrier and do 
if two agents will intersect with each other.
\subsubsection{Overlapping lanes}
One of the techniques involved in the focal mode of visual cognition includes estimating the future possible trajectories based on lane semantics. In other words, the interaction between agents can also be determined based on the lanes which are available in the future for each agent. Based on this information, we can track if there will be any possible intersection of lanes between two agents for them to have an interaction. Lanes for each agent can be found using lane tracking algorithms \cite{NGUYEN2018822}\cite{lane}. % such as feature-based tracking, model-based tracking, or optical flow. 
They can also be derived from map information assuming the data is accurate. Considering the map semantics to be accurate, we track the possible lanes of each agent. We have lane dividers spanning all the lanes on the roads. These dividers are defined by a set of token identifiers $\tau_{e_{0:N}}$ and are connected with each other by common end tokens $\tau_{e_0}$ and $\tau_{e_N}$. The direction of traffic in a particular lane divider can be obtained by extrapolating any of the tokens from $\tau_{e_{0:N}}$. Each agent can be associated with a lane divider $\tau_{d_i}$ based on its current position. To find if any two agents would meet in the future, we find if there's any intersection in the list of possible lanes of each agent $\tau_{l_i}$. In order to find $\tau_{l_i}$ for each agent $i$, we collect all lane dividers connected to $\tau_{d_i}$ by considering the dividers having $\tau_{e_N}$ as a common token. Consequently, we check for the intersection of the list of lane tokens between agents, $\alpha_{ij} = \mathbb{I}(\tau_{l_i} \cap \tau_{l_j})$.%, i.e. $\alpha_{ij}=0$ if there are no intersections, 1 otherwise. %In other words, the two agents are considered valid for the formation of cliques if they have at least one lane intersection between them. 


%\begin{figure}
%\begin{center}

%\includegraphics[width=1.0\linewidth]
%   {iterative\_cliques.pdf}
%\end{center}
%  \caption{SIMF utilizes both modes of perception to capture time-dependency in clique formation, by extracting semantic information at timestamps sampled from $0:T$. On the contrary, traditional methods can predict future poses at fixed, preset time steps using only the ambient mode of perception. (Note: Although the range of timestamps in each scene goes up to 20 s, we select scenes that align with the indicated timestamps for illustration.)
%  }
%\end{figure}
% \section{Dynamic ScePT}

% Given the optimized cliques with timestamp annotations, the CVAE model [scept] outputs joint trajectory predictions dynamically, thereby considering the interactions between nodes in different cliques.
% In order to obtain the active nodes and their relative states in each clique, the cliques are rearranged into batches having a list of active nodes with their corresponding relative state and edge histories. These batches are encoded into the feature vectors via LSTMs. The LSTM-encoded state edge history along with CNN encoded local map generates a discrete Gibbs distribution over the clique latent variable. The discrete latent variables are converted to one-hot encodings and each encoding is passed through the decoder. The desired trajectory is calculated by passing the one-hot encoding, state history encoding and map encoding. The desired trajectories and each latent encoding is passed to a policy network to obtain closed-loop trajectory predictions. The policy network is inspired by ScePT[scept]. 

\subsection{Multi-head Attention-based Global Encoding}
Once the cliques are formed, their local semantic maps and node state and edge history encodings are extracted and further processed through a multi-head attention mechanism to obtain global encodings. To understand the significance of this mechanism, let us consider the state-of-the-art architecture for trajectory prediction consisting of a CVAE model that encodes node history and local map information through a Gibbs distribution and passes through a policy network to generate trajectory predictions. Since agents follow a regular pattern and are highly correlated with each other, global map and state history information is crucial to capture the joint trajectory predictions. Specifically, in cases where agents are occluded from the ego vehicle (cross-traffic scenario), local map information and its associated state history information are not sufficient to capture the global information of the agents in the clique. Also, global context is required to capture the inter-clique information essential for minimizing the collision rate between cliques. Hence, we propose a multi-head attention mechanism by passing map encodings, node history, and edge history information as unique heads through a custom-designed transformer network. The model captures different dependencies and relationships between these heads, thereby enabling active interaction between ambient and focal modes of the vision system. Let the map $M^t_i$ of each agent $i$ be encoded via CNN to produce local encodings $M_{l_i}^t$, and similarly, state and edge histories $(h_{s_i}^t, h_{e_i}^t)$ be encoded into $(h_{s_{l_i}}^t, h_{e_{l_i}}^t )$ using LSTMs. The multi-head attention mechanism acts on these local features ($M_{l_i}^t,h_{s_{l_i}}^t, h_{e_{l_i}}^t $) to produce global features  ($M_{g_i}^t, h_{s_{g_i}}^t, h_{e_{g_i}}^t $) and appended with the local contextual information to form ($\vec M^t_i, \vec h_{s_i}^t, \vec h_{e_i}^t $).  
\begin{wrapfigure}{r}{0.57\textwidth}
    \vspace{-20pt}
{\footnotesize$$
w_{i, j}= \begin{cases}\operatorname{softmax}\left(FC\left(M_{l_i}^t,h_{s_{l_i}}^t, h_{e_{l_i}}^t \right)\right)_j, & \text { if } j\neq i \\ 0, & \text { if } j=i \end{cases}
$$}\vspace{-20pt}
\end{wrapfigure}
%\\ \operatorname{softmax}\left(FC\left(M_{l_i}^t, h_{s_{l_i}}^t,  h_{e_{l_i}}^t \right)\right)_{j-1}, & \text { if } j>i
 The multi-head attention mechanism structure consists of a fully connected layer $FC$ outputting the weights of an attention matrix $A$ populated on its non-diagonal elements.  In greater detail, $FC$ layer outputs a probability weight for all agents with respect to each and every other agent. The local features are then multiplied with the weights to get a global feature vector of the same dimension,
$
M_{g_i}^t=\sum w_{i, j} M_{l_i}^t
$ , 
$
h_{s_{g_i}}^t=\sum w_{i, j} h_{s_{l_i}}^t
$ and  $ h_{e_{g_i}}^t = \Sigma w_{i,j}  h_{e_{l_i}}^t$. 
%The two processes can be combined into a single efficient step using parallelization techniques by appending both features onto a single vector and populating the weights in an attention matrix, followed by aggregating all the individual agents' feature vectors onto a single list. Let  $\mu_l = [\vec M^0_l,\vec M^1_l,...,\vec M^{n-1}_l]$ where $n$ is the batch clique size, and $A = [w_{i,j}]_{n \times n}$. Thus the global features can be obtained by matrix multiplication of $A$ and $\mu_l$. The similar method is followed for $\vec h^i_{s_l}$ and $\vec h^i_{e_l}$. We assume a constant maximum  possible batch clique size and pad the vectors with -1 for $\vec M_l^i$ and 0 for $(\vec h_{s_l}^i, \vec h_{e_l}^i )$.
\begin{SCfigure}
\centering
\includegraphics[trim={0.3mm 3.0mm 0 0},clip,width=0.7\textwidth]
   {simf3.pdf}%\vspace{-0.2in}
   \caption{Overview of SIMF: It takes the graph, $G^t$, at each time step, $t$, and embeds it based on scene local and global semantics into $z^t$ to predict the future states of all interacting agents, $s_{pred}^{t:t+\eta}$. The process is repeated for all $t\in [0,T]$.} %Graph is generated with nodes as agents. Cliques are formed based on various techniques mentioned in Section 3.1, further passed on to get batch cliques containing node histories and local maps. Upon further processing, they are fed into the Multimodal attention mechanism to get global and local encodings. These encodings are converted to time-dependent latent variables and fed into a policy network to obtain closed-loop trajectory predictions.}
   \label{fig:simf}
\end{SCfigure}
\subsection{Temporal Latent Encodings}
In SIMF, we present a method that utilizes both modes of interaction to capture dynamic information by associating each node with a time-dependent latent variable. To achieve this, we encode the %output of the multi-head attention mechanism containing 
global information along with local information from both modes of visual perception into discrete latent variables $z^t_i = [z^t_1, z^t_2, z^t_3, ..., z^t_N]$ for each clique $C^t_{1:n}$ consisting of agents 1 to $N$. To elaborate, consider the temporal encoding  framework as a function $g$. We can write $z_i^{t} = g(\vec h_{s_i}^{t},\vec h_{e_i}^{t}, \vec M_i^t,  s_{{ref}_i}^{t:t+\eta})$ where $s_{ref}$ denotes the reference trajectories during training, and the former variables represent the local and global encodings concatenated together. Note that this is contrary to recent work that capture temporal information implicitly by incorporating influence from neighboring nodes using edge encoders \citep{salzmann2021trajectron}, which incorporates only the ambient mode of information. %Note that this is contrary to recent work that capture temporal information implicitly based on incorporates explicit methods of dynamic information by incorporating influence from neighboring nodes using edge encoders \citep{salzmann2021trajectron}, which incorporates only the ambient mode of information.   %To illustrate the importance of temporal dynamics in these latent variables, we consider a simple illustration scene shown in Figure 1. At the $5^{th}$ second ($10^{th}$ timestep),  agent A no longer interacts with agents E and F, since they move apart, and E is no longer in close proximity with agent A nor F is towards the direction of agent A. Similarly, new agents' interactions are captured between agents A, B and C based on the time-varying latent variable encodings.\\ %include explicit and implicit graphs
Consequently, the temporal information is ignored if the latent encodings are not iteratively captured based on timestamps, hurting prediction accuracy as indicated by our experimental results. The addition of the time-dependent latent variables has significantly improved the optimality in choice of agents, as illustrated in Figure 1 of supplementary material. %This can be seen in dense scenes, where the larger scene graphs are partitioned into small cliques. The addition of the time-dependent latent variables has significantly improved the optimality in choice of agents.
\subsection{Overall pipeline}
In SIMF, we generate a spatio-temporal graph $G^t$, with nodes representing agent and edges representing their interactions evaluated using the techniques mentioned in Section 3.1. The degree of interaction $\alpha$ is populated in an adjacency matrix to form optimal cliques $C^t$, i.e., sub-graphs of agents which have considerable interaction with each other. These cliques are further processed to obtain the active nodes and their relative states in each clique. Consequently, they are rearranged into a list of active nodes with their corresponding relative state and edge histories, i.e., the batches are populated with the aggregated information of all the nodes along with their corresponding clique details. Further, we generate local semantic maps \cite{caesar2020nuscenes} $M^t$, by rotating and cropping according to each agent's heading and type. These local maps are passed on to CNN to get local map encodings $M_l^t$,  while the state and edge histories are encoded into the feature vectors $(h_{s_l}^t, h_{e_l}^t)$ via LSTMs. These vectors are passed through the Multi-head attention-based mechanism to obtain global encodings  ($M_g^t, h_{s_g}^t, h_{e_g}^t $). The global encodings appended with the local encodings to form ($\vec M^t, \vec h_{s}^t, \vec h_{e}^t $) containing the complete information of the scene, further passed through the Gibbs discrete samplingo to produce latent variable encodings for each timestamp $0:T$ to capture the temporal information. These time-dependent latent variables $z^t$ are passed through the decoder containing the policy network to obtain the joint trajectory predictions $s_{pred}^{t:t+\eta}$, as illustrated in Figure \ref{fig:simf}. 
%  \section{Contributions}
% Two improvements: 
% \begin{itemize}
%     \item Choice of cliques based on the distance between them as a parameter is insufficient to determine optimal cliques. Used better mechanisms to find the cliques, as the quality of cliques affects the accuracy and computational efficiency of the network.
%     \item Dynamic cliques: If the cliques are formed based on the initial history states, there is no information update on the nodes during run-time. Formed cliques at an interval of time-steps to produce dynamic information in the choice of cliques.
% \end{itemize}
% 3 major modifications:
% \begin{itemize}
%     \item Using the semantics from map information:
%         Fine-tune the cliques by using only the nodes which are in possible interaction span. This can be obtained from semantics such as physical barriers in roads and possible lanes available for each node
%     \item Include the direction of motion between agent interactions to get a better estimate of cliques. For example, vehicles moving past each other in opposite directions need not be considered in the same clique, even though the distance between is small (beneath a threshold).
%     \item Modified the network to include cliques having an update of agents' dynamic information. 
% \end{itemize}


%===============================================================================

% \section{Citations}
% \label{sec:citations}

% 	Citations can be made using either \textbackslash citep\{\} or \textbackslash citet\{\}, depending from the appropriateness. To avoid the citation moving to the next line, it is often a good practice to replace the space before with a tilde (\~{}) character.

%===============================================================================
 \begin{wraptable}{r}{7.25cm}\vspace{-60pt}
                    \begin{tabular}{ccccc}
\toprule
Method  &@1s & @2s & @3s & @4s  \\
\midrule
S-LSTM [24]& 0.47 & - & 1.61 & - \\
CSP[25]      & 0.46 & - & 1.50 & - \\
CAR-Net [26] & 0.38 & - & 1.35 & - \\
SpAGNN [27] & 0.35 & -& 1.23 & - \\
Trajectron++ [19] & \textbf{0.07} & \textbf{0.45} & \textit{1.14 }& 2.20 \\
ScePT(Best-of-3) & 0.40 & 0.80 & 1.36 &\textit{ 2.14} \\
\midrule
SIMF(Best-of-3) & \textit{0.32 }& \textit{0.48} & \textbf{0.72}
 & \textbf{1.18 }\\
\bottomrule
\end{tabular}
\captionof{table}{Comparison of FDE (in meters) with Baselines. Ours perform better among most of the baselines, specifically in the latter timesteps of the prediction horizon. %\textbf{Bold} indicates the best value.
}\vspace{-0.35in}
\label{table:baseline}
\end{wraptable}
\section{Experimental Results}
\vspace{-10pt}
\subsection{Evaluation Metrics and Baselines} 
We use the standard evaluation metrics, also employed by the prior methods \cite{chen2022scept, salzmann2021trajectron, agentformer, laformer} naming Average Displacement Error (ADE) and Final Displacement Error (FDE). These metrics are based on $L2$ distance. The former is between the ground truth and predicted trajectories, whereas the latter is between the predicted and ground truth final agents' positions associated with the given time stamps. In addition, we also introduce a new metric, named Mean Average Count (mAC) which is inspired by a metric in \cite{waymo}. The mAC is determined as follows. We compute the trajectory prediction of the highest confidence by choosing the Best-of-N extension from the encoder output, i.e., we take the best from the N highest probability modes of the latent variable encodings taken at a time. Next, we check if the prediction is within lateral and longitudinal thresholds $\eta/3\mu(v_x)$ and $\eta/3\mu(v_y)$, respectively, where $\eta$ is time length and $\mu$ is a scaling factor computed based on the velocities $v_x$ and $v_y$. Further details on $\mu$ are provided in our supplementary material (Section 1.2.4). Finally, if the prediction lies within the thresholds, then it is considered a true positive; otherwise considered a false positive. The average of false positives is calculated per clique to estimate the accuracy of each clique. The overall mAC is calculated as the mean of all the cliques in the scene. For all the metrics, the lower is better, and bold/italic marks the best/second-best value. 
%Similar to previous works \cite{chen2022scept} \cite{salzmann2021trajectron} \cite{agentformer} \cite{laformer} , our method for trajectory forecasting is evaluated on Average Displacement Error (ADE) and Final Displacement Error (FDE) metrics.
% 1.\textbf{ Average Displacement Error (ADE):} It is defined as the mean $l_2$ distance between the ground truth and predicted trajectories. We have used the Best-of-N extension technique \cite{chen2022scept} by taking the N highest probability modes from the latent variable encodings. In order to enhance the metric, we have refined the ADE by excluding agents with zero ground truth predictions, thereby removing noise from the metrics.  \\
% 2.\textbf{ Final Displacement Error (FDE): } $l_2$ distance between the predicted final position and the ground truth final position at each prediction timestep $0:H$. We additionally check if the timesteps match with the ground truth for each node and discard predictions that do not have corresponding ground truth. \\
%Apart from these metrics, we propose a new metric for calculating the accuracy of trajectory predictions: \\
%\textbf{ Mean Average Count (mAC): }Inspired from \cite{waymo}, we compute the trajectory prediction of the highest confidence by choosing the Best-of-N extension from encoder output. We check if the prediction is within a lateral and longitudinal threshold $(T_{lat},T_{lon})$of the ground truth trajectory at each timestep $H$. The threshold is set empirically based on the magnitude of $H$ as $T_{lat} = H/3\mu(v_x), T_{lon} = 2H/3\mu(v_y)$, where $\mu$ is a scaling function according to $(v_x, v_y)$ defined as $\mu(v) = max(0,min(1,(v - v_L)/(v_H - v_L)))/2 + 0.5$ for each agent, where $v_H$ and $v_L$ are empirically set to 11 and 1.4 respectively. If the prediction lies within the threshold, then it is considered a true positive; otherwise considered a false positive. The average of false positives is calculated per clique to estimate the accuracy of each clique. The overall mAC is calculated as the mean of all the cliques in the scene. For all the evaluation metrics, the lower the better. 
\begin{SCfigure}
    \centering
    \includegraphics[width = 0.7\textwidth, height= 4cm]{gtandpred.png}
    \caption{Visual comparison of ground truth (left) and prediction trajectories (right) for a particular scene for both vehicles (blue rectangle) and pedestrians (pink circle). We can observe that each agent has multimodal predictions, with high-confidence predictions aligning closely with the ground truth paths}
    \label{fig:my_label}
\end{SCfigure}

% \begin{figure}
% \begin{minipage}[h!]{.35\textwidth} % not "0.5\textwidth"
%     \includegraphics[ height = 5cm]{visualization_gt.png}
%     \caption{Ground truth visualization}
%     \label{fig:test1}
% \end{minipage}%
% \hfill % <-- new
% \begin{minipage}[l]{.35\textwidth}
%     \includegraphics[ height = 5cm]{predic.png}
%     \caption{Prediction visualization}
%     \label{fig:test2}
% \end{minipage}
% \end{figure}
% \begin{figure}
% \begin{minipage}[h!]{.475\textwidth} % not "0.5\textwidth"
%     \includegraphics[width=\textwidth]{ade_graph.png}
%     \caption{ADE vs number of samples}
%     \label{fig:test3}
% \end{minipage}%
% \hfill % <-- new
% \begin{minipage}[h!]{.475\textwidth}
%     \includegraphics[width=\textwidth]{collision_rate.png}
%     \caption{Collision rate vs time in seconds}
%     \label{fig:test4}
% \end{minipage}
% \end{figure}

% \includepdf[pages={1-3},fitpaper]{cliquepruning.pdf}
\label{sec:result}
Our results are compared against these baselines: Social LSTM \cite{7780479} where each agent is modeled with LSTMs and followed by the multiagent “Social” pooling layer; Convolutional Social Pooling (CSP) \cite{Deo_2018} which explicitly considers a fixed number of trajectory classes thereby converting prediction to a classification task;
  \begin{wraptable}{r}{7.9cm}\vspace{-8pt}
 \begin{threeparttable}
    \begin{tabular}{ccccc}
   
\toprule
\multirow{ 2}{*}{\# samples}  &
\multicolumn{2}{c}{ADE} & \multicolumn{2}{c}{mAC}\\
\cmidrule{2-5}
 & ScePT  & Ours & ScePT  & Ours\\
\midrule
%  1 &
% 0.67 &\textbf{ 0} \\

 2 &
0.35/0.11 & \textbf{0.22/0.07} & 0.33  & \textbf{0} \\



 3 &
0.31/0.11 & \textbf{0.26/0.07} & 1.33  & \textbf{0} 


\\
 5 & 0.31/0.07 & \textbf{0.17/0.06} &
\textbf{0.33}  & 0.67
\\ 10 & 0.35/0.07 & \textbf{0.27/0.03}
& 2 &\textbf{ 1.67}
\\

\bottomrule
\end{tabular}
\captionof{table}{Comparison of ADE (in meters) and mAC (in percentage) with ScePT. SIMF performs better for all samples with respect to ADE, for most samples with respect to mAC.}\label{table:results}

\end{threeparttable}\vspace{-10pt}
\end{wraptable}
 CAR-Net \cite{sadeghian2018carnet} solves the prediction task by localizing the regions of interest for each agent using LSTMs-based visual attention; SpAGNN \cite{casas2019spatiallyaware} where raw LIDAR data and semantic maps are encoded using CNN followed by GNN to produce probabilistic trajectories; Trajectron++ \cite{salzmann2021trajectron} and ScePT \cite{chen2022scept} where the agents are represented as a spatiotemporal graph and the predictions are forecasted through a CVAE model. 
%  \begin{minipage}[t]{0.6\textwidth}%
%               \centering
%                 \label{Tabelle 3}


%     \end{minipage}%
%     %1 0 3 0 2
%     \begin{minipage}[t]{0.4\textwidth}%
%                 \centering
%                 \label{Tabelle 4
%                 }
%                    \begin{tabular}{ccc}
% \hline
% \# samples & ScePT & Ours \\
% \hline
% %  1 &
% % 0.4/ \textbf{0.14} &
% % \textbf{0.16}/ 0.17 \\

%  2 &
% 0.35/0.11 &
% \textbf{0.22/ 0.07}
% \\
%  3 &
% 0.31/0.11 &
% \textbf{0.26/0.07}


% \\
%  5 &
% 0.31/  0.07 &
% \textbf{0.17/ 0.06 }
% \\ 10 &
% 0.35/ 0.07 &
% \textbf{0.27/ 0.03 }



% \\
% %0.53, 0.15 & 
% %0.20, 0.06 & 
% %0.25, 0.03 &
% %0.24, 0.06 & 
% %0.17, 0.08 & 
% \hline
% \end{tabular}
% \captionof{table}{ADE results of Vehicles/Pedestrians in meters. \textbf{Bold} represents the best value.} 

% \end{minipage}%
\subsection{Results}
In this section, we present the comparison and ablation analysis on nuScenes dataset \cite{caesar2020nuscenes}. The ADE and FDE are represented in meters, and mAC is represented as the percentage of misses,  i.e., the percentage of the average number of nodes having predictions that do not lie within the expected threshold region. Furthermore, we evaluate with Best-of-N metric similar to the baselines to establish a fair comparison.
% \begin{wrapfigure}{r}
% {0.4\textwidth}{\vspace{-10pt}}\includegraphics[width=5cm,height=4cm]{collisionrate.png}
% \caption{Estimating relative interaction between two agents in a scene}\vspace{-10pt}
% \end{wrapfigure}
% \textbf{Qualitative Comparison:} The figure 4 shows trajectory predictions in comparison with their corresponding ground truth visualizations. The rectangles in blue represent Vehicles of different types (eg. cyclists, cars, trucks), while the pink circles represent pedestrians. It is evident that the predictions closely align with the ground truth for both vehicles and pedestrians. We can also observe that the predictions are multimodal in nature and hence the best possible prediction with a high confidence value (Best-of-N) is used for downstream planning of the ego-vehicle.
%SIMF on nuScenes \cite{caesar2020nuscenes} vehicles gives better results among almost all the baselines as seen in

\textbf{Comparison with baselines:}
In this section, we compare our approach, SIMF, against the given baselines. Table ~\ref{table:baseline} summarizes the FDE-based comparison of all methods. It can be seen that SIMF gives overall better results among almost all the baselines, especially in the latter timesteps of the prediction horizon. This clearly shows that the temporal dependencies in lateral encodings are vital in reducing the accumulated error from the previous timesteps' predictions. Although Trajectron++ \cite{salzmann2021trajectron} gives the best result in the short term, it is computationally expensive to execute as it does not form cliques. Among baselines, only the ScePT forms cliques; therefore, we center our remaining comparison analysis against the ScePT method. % and is  Further, we conduct experiments comparing our approach with ScePT, as it facilitates highlighting the importance of clique formation using focal mode. 
Table \ref{table:results} compares our method with ScePT using mAC and ADE. The results show that our model performs better than the ScePT. Furthermore, as we increase the number of samples taken from the Gibbs distribution, there is an increase in miss rate, possibly due to the impact of the tunable risk measure incorporated from ScePT that modifies the confidence level of predictions.  
 \begin{wrapfigure}{r}
 {0.53\textwidth}{\vspace{-15pt}}\includegraphics[width=7.5cm]{collision.png}
\caption{SIMF demonstrates a reduced collision rate compared to ScePT, and the incorporation of an attention mechanism further decreases the collision rate, resulting in improved scene consistency.}\vspace{-20pt}
\label{fig:collision}
 \end{wrapfigure}
Specifically, there are high-confidence predictions that are different from the ground truth when increasing the number of sample encodings taken at a time. In terms of collision rate, SIMF shows significantly enhanced scene consistency than ScePT, as seen from the graphs in Figure ~\ref{fig:collision}, demonstrating a considerable decrease in collision rate, specifically with the use of the attention mechanism. Apart from the quantitative study, we qualitatively compare trajectory predictions in comparison with their corresponding ground truth visualizations in Figure ~\ref{fig:my_label}. %The rectangles in blue represent Vehicles of different types, while the pink circles represent pedestrians. 
It is evident that the predictions closely align with the ground truth for both vehicles and pedestrians. 

\textbf{Ablation study:} 
To develop an understanding of the individual components of SIMF, we conduct an ablation study having two parts. First, we ablate the model architecture on semantic aware clique formation, temporal encodings, and  attention mechanism. Secondly, we investigate the impact of semantic cliques by ablating the various techniques of Section 3.1, including direction, barrier detection, and overlapping lanes. We utilize ScePT for the ablation study since it quantifies the degree of interaction solely using the Euclidean distance parameter. Table ~\ref{table:ablation} summarizes our results, columns 2-5 and 5-9 highlight the first and second parts of our ablation study. It can be observed that the attention mechanism has the greatest impact on the accuracy, whereas the temporal factor has more impact on the later timesteps. From the second part, we can ascertain that the direction factor has the greatest impact among all the techniques, owing to its intuitive reasoning of relative velocities. 
%To develop an understanding of which components of semantics-aware clique formation (Section 4.1) influence performance, a comprehensive ablation study is performed in Table 4. We can see that the interaction detection using the direction parameter has the maximum impact on the performance compared to other parameters. Without the parameter, the model performs almost similar to ScePT, when comparing FDE values at all prediction horizons. 
%Table 6 corresponds to understanding the importance of different modules (Section 4.1, 4.2, 4.3) and its impact on the FDE. 


    %1 0 3 0 2
    % \begin{figure}
        


    %             \centering
    %             \label{Tabelle 2}
                 
  
%     \begin{minipage}[t]{0.5\textwidth}%
%               \centering
%                 \label{Tabelle 1
%                 }
%                \setlength{\tabcolsep}{2.5pt}
%  \begin{threeparttable}
%                     \begin{tabular}{ccccc}
% \hline
% Method  &@1s & @2s & @3s & @4s  \\
% \hline
% S-LSTM [24]& 0.47 & - & 1.61 & - \\
% CSP[25]      & 0.46 & - & 1.50 & - \\
% CAR-Net [26] & 0.38 & - & 1.35 & - \\
% SpAGNN [27] & 0.35 & -& 1.23 & - \\
% Trajectron++ [19] & \textbf{0.07} & \textbf{0.45} & \textit{1.14 }& 2.20 \\
% ScePT(Best-of-3) & 0.40 & 0.80 & 1.36 &\textit{ 2.14} \\
% \hline
% SIMF(Best-of-3) & \textit{0.32 }& \textit{0.48} & \textbf{0.72}
%  & \textbf{1.18 }\\
% \hline
% \end{tabular}
%  \captionof{table}{FDE comparison with Baselines. \textbf{Bold}/\textit{italic} indicates the best/second-best value.}
%  \end{threeparttable}
%     \end{minipage}\par

%    \end{figure}

  
%===============================================================================
\begin{table*}
\centering
  \begin{tabular}{ccccccccc}
   
\toprule
time stamp  &SIMF & w/o S & w/o T &  w/o A & w/o B & w/o O & w/o D & ScePT\\
\midrule
@0.5s & \textbf{0.19} & 0.28 & 0.22 & 0.33 & 0.19 & 0.20 & 0.31 & 0.33 \\
 @1.5s  &\textbf{ 0.41} & 0.50 & 0.46 & 0.55 & 0.43 & 0.44 & 0.50 & 0.54 \\

  @2.5s&\textbf{0.56}& 0.69&0.62& 0.76 & 0.60 & 0.61 & 0.68 & 0.76\\
@3.5s & \textbf{0.93} & 1.08 & 0.99 & 1.21 & 0.92 & 0.92 & 0.76 & 1.22\\
\bottomrule
\end{tabular}
\captionof{table}{Ablation study of techniques presented in SIMF according to FDE in meters. Legend: S- Semantic Cliques, T - Temporal encodings, A - Attn mechanism, B - Barrier Detection, O - Overlapping lanes, D - Interaction using direction. Attention mechanism (A) has the highest influence on SIMF, followed by the interaction using direction parameter (D) in semantic-aware clique formation techniques.}
\vspace{-0.2in}\label{table:ablation}\end{table*}
\section{Limitations and Future Work}\vspace{-10pt}
%In this paper, we propose SIMF, a CVAE-based model that generates multi-modal joint motion predictions by capturing focal and ambient modes of visual perception.
% It generates a graph with agents as nodes and breaks it into cliques based on semantic and position information of agents. Further, it converts the clique information into encodings containing local node information and passes it through an attention mechanism to obtain global context, thereby capturing both modes of perception. Finally, it generates time-dependent lateral encodings to produce closed-loop trajectory predictions via a policy network.
%The results show that SIMF gives better results among all the baseline models, achieving improved collision rates and scene consistency. % motion prediction accuracy along with improved collision rates and scene consistency. %We also propose a new metric for evaluating the prediction accuracy of models, owing to the need for simplicity in evaluating the predicted trajectories.
One of the limiting factors of our approach is that we assume road semantics information is available and accurate, which might not be the case in real-world settings. Therefore, our future work includes working on an end-to-end model using only camera inputs to provide solutions to vehicles that lack expensive onboard sensors for feature extraction. Secondly, our method implicitly captures inter-clique interaction, which with explicit modeling, may lead to further improvements in our performance. Therefore, another of our future works is explicitly capturing inter-clique interactions that may be achieved using reparameterization techniques to backpropagate through the time-dependent latent variables sampled during training for temporal latent encodings. %Finally, SIMP can be improved in terms of computational complexity by implementing efficient coding and parallelization techniques. 
\label{sec:conclusion}

%===============================================================================

%\clearpage
% The acknowledgments are automatically included only in the final and preprint versions of the paper.
% \acknowledgments{If a paper is accepted, the final camera-ready version will (and probably should) include acknowledgments. All acknowledgments go at the end of the paper, including thanks to reviewers who gave useful comments, to colleagues who contributed to the ideas, and to funding agencies and corporate sponsors that provided financial support.}

%===============================================================================

% no \bibliographystyle is required, since the corl style is automatically used.
\bibliography{example}  % .bib
\vspace{80pt}
\section{APPENDIX}
\subsection{Algorithmic Procedures and Coding Implementation:} Algorithm ~\ref{alg:two} provide our pseudocode for semantic-aware clique formation. Further details and overall implementation code will be made publically available with the final paper release. 
\RestyleAlgo{ruled}
\begin{algorithm}[hbt!]
\caption{Semantic-aware clique formation }\label{alg:two}
\KwData {D $\rightarrow$ $D_{train}$, $D_{validation}$}

\For {$V$ in $D_{train}$}
{$\tilde{V}_{0:T}$ $\leftarrow$ PresenceTable($V_{s_{0:T}}$)\Comment list of active nodes for $T$ timesteps\\} 
\For {t in T}  
{
%$x_{{traj}_t}$ $\leftarrow$ PropagateTraj($h_{s}^{t-1}$, $dyn_n$, $con_n$) \\
$d_{ij}\leftarrow \min ||\phi_{a_i}^t(s_i), \phi_{a_j}^t(s_j)||_2$ \Comment for all ij pairs of active nodes $\tilde V_{0:T}$\\
$\delta_{ij}\leftarrow ||(p_o - p_i) + t (v_o - v_i)||_2 $ \Comment Compute relative velocity vector\\
\If {$\delta_{ij} < D$ and $i \neq j$ and $d_{ij} < d_0$}
{

$m_{AC},m_{AD},m_{BC},m_{BD} \leftarrow $CalculateSlope$(\tau_{A\rightarrow{B}}, \tau_{C \rightarrow{D}})$\\
$barrier \leftarrow$ ($m_{AC}<m_{AD}) \oplus (m_{BC} < m_{BD}$) \Comment Check if CD intersects AB\\
\If {$s_{v_i}$and $s_{v_j}$ and not $barrier$} 
{
$\tau_{l_i}$, $\tau_{l_j}$ $\leftarrow$ GetNextLanes($\tau_{d_i},\tau_{d_j}$) \Comment Collect all possible future lanes $\forall \tau_{d_i},\tau_{d_j}$\\
\If {$\tau_{l_i} \cap \tau_{l_j}$}
{
$\alpha_{ij} = d_0/d_{ij}$ \Comment{Update $\alpha_{ij}$}
}
}
}
$C^t_{1:n} \leftarrow$ connectedGraph($\alpha$) \Comment{connected components of  $\alpha$}\\

return $C_{1:n}^t$

}

\end{algorithm}

\subsection{MODEL ARCHITECTURE}

Our model consists of a CVAE framework encapsulating two parts: an encoder containing a Multi-head attention network and a decoder containing a policy network. The input for our framework includes the output from the Semantic-aware clique formation module $C^t$, which we convert into time-dependent latent variables $z^t$  to produce $s_{pred}^{t:t+\eta}$. The model components are described as follows:

Each agent in the scene is considered to have interaction with atleast one other agent, and hence all the agents are involved in clique formation. Also, each agent can be a part of multiple cliques, and hence can have more than $N$ modal predictions. We take the Best-of-N for metrics evaluation, if the number of prediction are greater than $N$. The ego-vehicle becomes part of the cliques similar to any other agent in the scene during prediction, followed by a planning module which takes in these predictions and implements the ego-motion forecasting.
\subsubsection{Encoder} 
The encoder models the neighboring agents' influence for each agent in three steps. First, we capture the local information of each agent by passing it through LSTMs and CNN. Each agent's node state and edge history information are passed through LSTM networks to produce local encodings $(h_{s_l}^t, h_{e_l}^t)$. The agent-agent interaction are classified into 3 types: Vehicle-Vehicle, Vehicle-Pedestrian, Pedestrian-Pedestrian. Each interaction type share a common set of weights for the LSTM.  Specifically, each node is classified into a Vehicle or pedestrian and is aggregated into two separate lists before passing through the LSTM.  Further, the semantic maps of each agent are converted to map encodings $M_l^t$ using a CNN followed by an FC layer, with a leaky ReLU activation function. Finally, the node and edge histories are concatenated together to form $(M_l^t, h_{s_l}^t, h_{e_l}^t)$. \\
\\
Secondly, we propose a self-attention
distillation (SAD) module \cite{hou2019learning} to aggregate contextual information $(M_g^t, h_{s_g}^t, h_{e_g}^t)$ using custom-designed transformers. The module establishes active interaction between the map encodings containing focal mode information, node and edge histories containing ambient mode information, simultaneously capturing global interaction between the nodes of different cliques (inter-clique information).  The module allows for a lightweight structure containing simple FCs and softmax modules, thus achieving highly accurate global encodings while maintaining computational efficiency.  Specifically, we pass  $(M_l^t, h_{s_l}^t, h_{e_l}^t)$ through an FC layer, processing each local feature vector  to populate an attention matrix $A$ containing probability weights $w_{i,j}$, as explained in Section 3.2. \\
\\
Thirdly, these global and local encodings $(\vec M^t, \vec h_{s}^t, \vec h_{e}^t)$ along with reference trajectories $s_{ref}^{t:t+\eta}$ during training are passed through a temporal encoding network $g(.)$ to represent the encodings as a Gibbs distribution $P(z^t)$ by introducing a discrete latent variable 
$z^t_i = [z^t_1, z^t_2, z^t_3, ..., z^t_N]$ for each clique $C^t_{1:n}$ consisting of agents 1 to $N$ (Section 3.3). $g(.)$ consists of feed-forward neural networks, mapping each of $(\vec M_i^t, \vec h_{s_i}^t, \vec h_{e_i}^t)$ to its corresponding $z_i^t$. The log-likelihood of $P(z^t)$ can be computed by constructing a factor graph, containing $z_i^t$ as variable nodes and network weights as edges. Specifically, each possible valuation of $z_i^t$ is obtained from the graph by summing all the factor nodes of $z_i^t$. 
\begin{figure}
    \centering
    \includegraphics[width = 1.0\textwidth]{temporal.pdf}
    \caption{An illustration of the cliques changing over time. Fig \ref{fig:temporal}(3) from left shows that agents E and F no longer in interaction with the ego-vehicle A are removed from the clique at the $10^{th}$ timestep (5 seconds), whereas C and B are added into the clique which was not included earlier. This shows the impact of temporal information in cliques and, thus, the importance of time dependency imbibed in latent variables.}
    \label{fig:temporal}
\end{figure}
\subsubsection{Decoder}
SIMF uses a policy network to model each agent as a motion planner and learn a policy to emulate its own future path to output trajectory predictions. The policy network takes in the encodings  $(\vec M^t, \vec h_{s}^t, \vec h_{e}^t)$, reference trajectories $s_{ref}^{t:t+\eta}$, latent variable $z^{0:T}$ sampled from the Gibbs distribution, and outputs $s_{pred}^{t:t+\eta}$.  In particular, the policy network uses a feed-forward action model $\Omega$ followed by a Dynamics function pertaining to each agent type to learn a policy that outputs predictions. The Dynamics function is assumed to be a differentiable function of the state and control input, which encapsulates most types of real-world agents. The pedestrians are modeled as single integrators, and each wheeled vehicle is modeled as a dynamically extended unicycle, instead of learning the parameters online for each type of vehicle. For each timestep $t$, the latent variables are fed into the policy network to produce a corresponding $s_{pred}^t$, thus facilitating the policy to capture the dynamic information of latent variables. Specifically, at a timestep $t$, the inputs  $(\vec M^t, \vec h_{s}^t, \vec h_{e}^t, z^t, s_{ref}^t)$ are passed onto $\Omega$ to produce control actions, fed into the Dynamics function to output the agent's state $s_{pred}^t$.The state prediction is then compared with $s_{ref}$ to compute a tracking error $\Delta s$, which is fed back to the policy network  to produce closed-loop predictions. The process is repeated for $\eta$ steps to obtain $s_{pred}^{t:t+\eta}$. \\
\\
We use dynamically extended unicycle model from \cite{salzmann2021trajectron}, which uses acceleration $a$ and heading rate $\omega$ as input $\mathbf{u} = [\omega,a]^T$ to the model, instead of the unicycle generic model which uses velocity and heading rate. At each timestep, and for a specific latent value $z^t$, SIMF produces a 2D Gaussian distribution over control actions $\mathcal{N}\left(\mu_{\mathbf{u}}, \Sigma_{\mathbf{u}}\right)$. The controls $\mu_{\mathbf{u}}$ and uncertainties $\Sigma_{\mathbf{u}}$ are then integrated through the dynamics $s_{pred}^{t+1} = f(s_{pred}^t, \mathbf{u}^t)$ to obtain the future state distribution using the extended Kalman filter techniques from \cite{thrun2000probabilistic}. 
\subsubsection{Average and Final Displacement Error
Evaluation}
1.\textbf{ Average Displacement Error (ADE):} It is defined as the mean L2 distance between the ground truth and its corresponding predicted trajectories. We have used the Best-of-N extension technique \cite{chen2022scept} by taking the N highest probability modes from the latent variable encodings. In order to enhance the metric, we have refined the ADE by excluding agents with zero ground truth predictions, thereby removing noise from the metrics.  \\
2.\textbf{ Final Displacement Error (FDE): } At each prediction timestep from $0:\eta$, the L2 distance between the predicted final position and the ground truth final position is calculated. We additionally check if the timesteps match with the ground truth for each node and discard predictions that do not have corresponding ground truth.
\subsubsection{mAC in detail}
When averaging the measurements in the case of ADE and FDE, the inherent probabilistic details imbibed from generative methods are lost. Hence we use mAC to quantify the miss rate of individual predictions, in order to compare the model and baselines with a practical outlook. 
The threshold for latitude and longitude is set based on the magnitude of $\eta$ as $T_{l a t}=\eta / 3 \mu\left(v_x\right), T_{l o n}=2 \eta/ 3 \mu\left(v_y\right)$, where $\mu$ is a scaling function according to $\left(v_x, v_y\right)$ defined as $\mu(v)=\max \left(0, \min \left(1,\left(v-v_L\right) /\left(v_H-v_L\right)\right)\right) / 2+0.5$ for each agent, where $v_H$ and $v_L$ are empirically set to 11 and 1.4 respectively. 
\subsubsection{Nuscenes dataset}
We conduct our experiments on the nuScenes dataset \cite{caesar2020nuscenes}, which is a comprehensive dataset for autonomous driving consisting of 1000 scenes located in Boston and Singapore, providing a large-scale environment for evaluation. Each scene is annotated at a frequency of 2 Hz, starting from time t = 0.5s, and spans a duration of 20 seconds. The dataset encompasses 23 different semantic object classes, and the high-definition semantic maps consist of 11 annotated layers, which are utilized for focal mode of perception.


\subsubsection{TRAINING AND IMPLEMENTATION DETAILS}
SIMF was implemented in PyTorch on a server running Ubuntu 18.04 docker containing 3.50GHz × 8 Intel Core i9 processor, 32 GB RAM, and GeForce RTX 3090 GPU. We trained the model for 100 epochs, each of batch size 256 (12 hours) on the nuScenes dataset, containing both Vehicle and Pedestrian information, from multi-modal fully autonomous vehicle sensor inputs: 6 cameras, 5 radars and 1 lidar.

\textbf{CVAE:} We use the architecture proposed by \cite{chen2022scept} with the encoder modified to output time-dependent latent encodings.  The model contains 4 major components; CNN for semantic map encoding, Gibbs distribution for encoder, Multi-head attention mechanism (Section 3.2) and policy network for decoder.

The map encoder consists of a CNN which takes in local maps $M^t \ \epsilon \ \mathbb{R}^{100 \times 100 \times 3} $, feeds it through 4 convolutional layers having channels of size $\{10,20,10,1\}$ with filters $\{5,5,5,3\}$ and corresponding strides $\{2,2,2,1\}$  followed by an FC layer of 32 hidden dimensions to produce local map encodings. %dropout?

The multi-head attention mechanism comprises of FC layers of 299 hidden dimensions to update the non-diagonal elements of the attention matrix of size $300 \times 300$. Since each batch has a different number of nodes froming the cliques, the encodings are padded with $-1$ for all heads.  Further, the local map, state and edge encodings for each node represented as a 32 dimensional vector respectively act as unique heads and gather global encodings of similar size to form concatenated 128-dimensional vector encoding.  

The encoder uses Gibbs distribution to capture the temporal ambient-focal information to produce discrete time-dependent latent variables. It contains two simple FC layers each of 64 hidden dimensions followed by a  LSTM with 32 hidden dimensions, fed into a softmax layer to represent the state and edge encodings of 128 dimension as $z^t$ for each timestep $t$ in $0:T$ of size considered to be $N=4$ as the max clique size. 

The decoder contains a 128-dimensional Gated Recurrent Unit (GRU), followed by LSTMs, fed by the temporal latent variables and concatenated 128-dimensional vector encoding. The output of the GRUs are further passed on to action net $\Omega$ which consists of 2 simplelinear FC layers of hidden layers $[128,64]$ to produce 2 outputs defining the control actions (acceleration and steering rate). These actions are further passed on to the Dynamics function.


%\bibliography{example} 
\end{document}
